\chapter{Create the Bot with BotFather}

\begin{goalbox}
Create a Telegram bot account and make its token available safely to Python.
\end{goalbox}

\section{Ask BotFather for a bot}

\begin{enumerate}
\item Open Telegram and search for the verified \texttt{@BotFather} account.
\item Start the chat and send \texttt{/newbot}.
\item Enter a display name, for example \texttt{Alireza Tutorial Bot}.
\item Enter a unique username ending in \texttt{bot}, for example
\texttt{alireza\_tutorial\_123\_bot}. If it is taken, choose another.
\item BotFather returns a token. Treat it like a password.
\end{enumerate}

\begin{warningbox}
Do not put the real token in source code, screenshots, chat messages, or Git. If
it is exposed, send \texttt{/revoke} to BotFather and generate a new token.
\end{warningbox}

\section{Set the token for one terminal session}

Replace the placeholder with the token BotFather gave you.

\textbf{macOS or Linux:}
\begin{lstlisting}[language=bash,numbers=none]
export BOT_TOKEN="paste_your_token_here"
\end{lstlisting}

\textbf{Windows PowerShell:}
\begin{lstlisting}[language={},numbers=none]
$env:BOT_TOKEN="paste_your_token_here"
\end{lstlisting}

The variable disappears when that terminal closes, which is safe and simple for
a first project. You must set it again in each new terminal session.

\begin{checkpointbox}
Check only whether the variable exists; do not print its secret value.
\begin{lstlisting}[language=bash,numbers=none]
python -c "import os; print('token set' if os.getenv('BOT_TOKEN') else 'missing')"
\end{lstlisting}
The expected result is \texttt{token set}.
\end{checkpointbox}
